\section{Tindak Lanjut dan Topik Lanjutan}
\label{sec:ch17-lanjut}

\subsection{Setelah Front-End}

Buku ini berfokus pada front-end. Topik lanjutan yang dapat dipelajari:
\begin{itemize}[leftmargin=*]
  \item \textbf{representasi antara}: SSA, three-address code, LLVM IR;
  \item \textbf{optimasi}: constant folding, dead code elimination,
        loop optimization, data-flow analysis;
  \item \textbf{pembangkitan kode}: instruction selection, register
        allocation, calling convention;
  \item \textbf{sistem tipe}: polimorfisme, inferensi tipe,
        dependent types;
  \item \textbf{analisis program}: static analysis, abstract
        interpretation, model checking;
  \item \textbf{pengujian compiler}: fuzzing, differential testing,
        metamorphic testing;
  \item \textbf{bahasa domain khusus}: desain dan implementasi DSL.
\end{itemize}

\subsection{Jalur Belajar}

\begin{table}[H]
\centering
\caption{Jalur belajar lanjutan}
\label{tab:ch17-jalur}
\begin{tabularx}{\textwidth}{@{}l X l@{}}
\toprule
\textbf{Minat} & \textbf{Topik} & \textbf{Sumber} \\
\midrule
Compiler engineering & LLVM, codegen, optimasi & \code{llvm.org} \\
Bahasa pemrograman & sistem tipe, semantik formal & buku TAPL \\
Analisis program & static analysis, abstract interpretation &
riset dan alat open source \\
Pengujian & fuzzing, differential testing & \code{oss-fuzz} \\
DSL & parser generator, interpreter & ANTLR, Xtext \\
\bottomrule
\end{tabularx}
\end{table}

\subsection{Karier}

Kompetensi front-end berguna pada:
\begin{itemize}[leftmargin=*]
  \item pengembangan compiler dan alat bahasa;
  \item static analysis dan linter;
  \item query optimizer dan mesin data;
  \item pengembangan DSL dan konfigurasi;
  \item keamanan aplikasi (analisis kode);
  \item AI untuk kode (parsing, transpiling, verifikasi).
\end{itemize}

\subsection{Saran Akhir}

Terus berlatih dengan proyek nyata. Baca kode compiler open source,
tetapi mulai dari modul kecil. Dokumentasikan keputusan desain dan
uji secara sistematis. Kemampuan membaca spesifikasi, memodelkan
bahasa, dan memberikan diagnostik yang baik akan berguna jauh di
luar mata kuliah ini.

\begin{infobox}
Compiler adalah sistem besar yang dibangun dari komponen kecil yang
saling terhubung. Keterampilan yang Anda latih di sini---membaca
spesifikasi, menulis parser, memeriksa tipe, dan menguji---berlaku
pada banyak sistem perangkat lunak lain.
\end{infobox}
